\subsection{3.2.1. 持久化 KV 缓存管理}\label{persistent-kv-cache-management}

在相同负载下，V4.1 将持久化 KV 缓存的占用降至 V4 的 \(1 / 8\)。这一缩减由相乘的两个因素共同实现：持久化 KV 缓存不再存储 SWA KV，使其规模几乎减半；缓存中保留的全局 KV 又经由架构与精度优化被进一步压缩至 V4 的 1/4。

在 V4 部署中，SWA KV 几乎占去持久化 KV 缓存容量的一半。在此缓存中，全局 KV 与 SWA KV 由 LRU 淘汰策略独立管理。全局 KV 被完整存储，命中后即复用完整前缀。相比之下，SWA KV 被缓存于两个特定位置——提示末尾与输出末尾——以支持重生成与多轮会话；命中时计算可从该缓存位置继续。为维持较高的命中率，我们在 SSD 上配置了足够大的持久化 KV 缓存，使典型工作负载下两类 KV 均能常驻超过 72 小时。尽管仅在指定位置保留 \(n _ { \mathrm { W i n } }\) 个 KV 条目，未压缩的 SWA KV 缓存仍带来可观的存储开销，尤其是在轮次较短的多轮对话中。

持久化存储 SWA KV 既代价高昂又收效甚微，因为其访问模式与持久化 KV 缓存的长效保留策略不匹配。与表现出长尾复用的全局 KV 不同，SWA KV 只在活动会话内部一个狭窄、以分钟计的窗口内被复用，一旦会话结束或下一轮开始即告失效。V4 技术报告提出了 Zero SWA Caching，通过重算缺失的 SWA KV 来规避存储开销。然而，精确恢复需要对 \(L \times n _ { \mathrm { w i n } }\) 个 token 执行一次完整的前向传播，这一成本在生产部署中被证明是不可承受的。

因此，V4.1 对持久化 KV 缓存的管理方式修订如下：

\begin{enumerate}
\def\labelenumi{\arabic{enumi}.}
\item
  SWA KV 不再缓存于持久化 KV 缓存，而是存入由每台机器 10\% 的主机 DRAM 划拨形成的分布式内存池。尽管该池在总容量上远小得多，但其较短的 TTL（仅有几分钟）允许过期条目被立即回收用于新会话；在真实工作负载下，这种高周转率足以服务绝大多数并发活动会话。全局 KV 仍保留在持久化 KV 缓存中，保证其生命周期至少为 72 小时。
\item
  逐出 SWA KV 不可避免地会造成未命中，而一种轻量级兜底机制——编码器 SWA 有界回放（详见 3.2.2 节）——使这类未命中仍可承受。对于那些命中全局 KV 但未命中 SWA KV 的必然却罕见的请求，它只对 \(n _ { \mathrm { w i n } }\) 个 token 进行重算，而非执行一次完整的 \(L \times n _ { \mathrm { w i n } }\) token 前向传播，从而恢复缺失状态。这种有界回放是整套设计的基础：它将灾难性的未命中转化为优雅而廉价的性能降级，进而为移除持久化 KV 缓存中的 SWA KV 提供了依据。
\end{enumerate}



\subsection{3.2.2. SWA 有界回放}\label{swa-bounded-replay}

由于 SWA 依赖关系在各层间累积，对 $L$ 层 SWA KV 的精确重建需要回放 \(L \times n _ { \mathrm { w i n } }\) 个 token。SWA 有界回放只回放最近的 \(n _ { \mathrm { { w i n } } }\) 个 token，并将 SWA 截断至回放片段，从而接受近似状态：对于从位置 $s$ 开始的回放，位置 $i$ 的查询会注意 $[\max(s, i - W + 1),\ i]$ 范围内的 SWA 键。

编码器 SWA 有界回放。编码器 SWA 有界回放使前缀缓存只依赖全局 KV，从而可将 SWA KV 从持久化 KV 缓存中移除。

当编码器 SWA KV 缺失时，我们回放已缓存前缀的最后 \(n _ { \mathrm { W i n } }\) 个 token，并将其与未缓存的后缀一同处理。被回放的 token 仅重新生成 SWA KV，并复用已缓存的全局 KV，既不重算也不覆盖；而未缓存的后缀则同时生成全局 KV 与 SWA KV。

按设计，回放得到的前缀状态是近似的，因此为未缓存后缀计算出的全局 KV 与 SWA KV 依赖于缓存命中位置，在不同位置之间并不数学一致。令人鼓舞的是，我们的实验证据表明，这种有界回放策略对响应质量的损害可以忽略不计。

解码器 SWA 有界回放。解码器 SWA 有界回放将解码器前向传播限定在 \(n _ { \mathrm { w i n } }\) 个 token 之内，使预填充总计算量几乎减半。

在 CED 下，解码器全局 KV 由编码器最终隐藏状态投影而来。在编码器处结束预填充的唯一障碍是解码器 SWA KV，它由每个解码器层自身的隐藏状态生成，并为最初的若干解码步骤所需。由于我们从不缓存解码器 SWA KV，精确重建它需要让 \(\frac { L } { 2 }\) 个解码器层处理最后 \(\frac { L } { 2 } \times n _ { \mathrm { w i n } }\) 个提示 token，当较短的未缓存后缀紧跟较长的已缓存前缀时，这一过程代价高昂。因此，我们也把有界回放策略应用于该场景：在每次预填充时，我们回放提示的最后 \(n _ { \mathrm { W i n } }\) 个 token，在相同的 SWA 截断下将其编码器输出经过解码器各层，并把由此得到的解码器 SWA KV 仅用于解码，而不用作前缀缓存。

按设计，重建出的解码器 SWA KV 与完整解码器前向传播所得在数学上并不等价。同时我们发现，这一策略对响应质量仅有可忽略的影响。为更加稳妥，我们还在后训练阶段模拟同样的回放，实现面向训练的适配。

\subsection{4. 预训练}\label{pre-training}

\subsection{4.1. 数据构建}\label{data-construction}

文本数据策展。为追求更高的智能水平，我们不再局限于小规模数据实验所反映的通用、样本级质量，而是更关注那些能带来独特信息增益的多元语料之间的整体互动。我们采用更系统、更规范的数据构建流水线，以提升数据质量并优化数据配比。具体而言，在更全面的评测基础上，我们精心设计了一个横跨模型参数与训练数据的缩放阶梯（scaling ladder），用以指导大规模训练。我们过滤掉信息增益有限的模型生成内容，包括能力较弱的模型产生的输出以及低质量的机器翻译文本。我们认为此类内容属于隐式重复，因为它在很大程度上只是对既有信息的改述，并可能在长期训练中产生负面影响。我们还探索了模型在环（model-in-the-loop）的数据迭代方法，作为未来大规模合成数据的基础。此外，我们引入更多领域专家，构建细粒度的数据质量评测维度。与前一版本相比，新语料纳入了更多来自新发布的开源仓库、代码提交、代码库与新兴框架的近期代码，从而覆盖更广泛的编程语言，更好地反映当代真实世界的软件工程场景。



多模态数据策展。我们的多模态预训练数据集主要包括三类数据：图文对、交错图文数据与领域专属数据。鉴于原始网页数据天然蕴含丰富的多模态知识，我们没有进行大规模数据合成，而是优先对数据进行清洗并以其原生形式加以利用，以便在预训练阶段实现最直接、可扩展的视觉知识压缩。在最初的数据采集过程中，我们发现抓取系统过度偏向以文本为中心的网页内容，因此我们以 Common Crawl 为源对其重新引导，以提升对多模态来源的覆盖。对于图文对，我们从网页中提取图像及其关联的 alt 文本，通过图文相关性阈值进行过滤，并基于图像语义进行去重。对于交错数据，我们主要从网页与 PDF 构建这一子集。处理大规模多模态语料通常比处理纯文本语料产生更高的 CPU 与磁盘存储开销，这促使我们将交错数据的构建组织为逐渐增耗的阶段。在图像检索之前，我们先施加启发式与统计过滤、去重及质量模型，以筛选出高价值的文档。通过筛选的文档随后被组装成交错图文序列，此时再次以图像感知的方式进行过滤与去重。最后，我们使用 SmolVLM（Marafioti 等人，2025）对图文内容进行严格的质量评分，从而提取高质量的交错数据。该过程中被过滤掉的文档有一部分经由筛选与重组，被回收为额外的图文对。为弥补从网络采集数据的固有局限，我们还纳入领域专属数据集，以增强模型在细粒度视觉感知（例如视觉定位与指认）、光学字符识别（OCR）以及长尾知识获取方面的能力。我们还收集了大量图像—代码对与计算机使用轨迹，以提升多模态智能体理解能力。

数据整合与去重。由于我们的纯文本数据与多模态数据经由不同的流水线处理，我们把最终训练语料构建为两类数据源的并集。对于重叠样本，我们用多模态版本替换纯文本版本，并采用两种配置中较大的轮数。完成替换后，最终语料的纯文本与多模态数据按 7:1 的 token 比例构成。我们通过联合预取与分配训练样本，将预训练与上下文扩展期间的样本重叠降至最低。超长文档在混料之前被确定性预切分，以确保训练 token 在数据分片与训练步骤之间均匀分布。我们还进一步增强了最佳适配打包（best-fit packing）算法，将填充率降至至多 10−4。